% Latin-script Interslavic source component: bibliography.
% Audited editorial provenance: ../00_SESSION_INDEPENDENT_STATE_v019.json and ../NORMALIZATION_DECISIONS_v019.jsonl.
\documentclass[11pt]{article}
\usepackage{fontspec}
\IfFontExistsTF{Times New Roman}{\setmainfont{Times New Roman}}{\setmainfont{TeX Gyre Termes}}
\IfFontExistsTF{Arial}{\setsansfont{Arial}}{\setsansfont{TeX Gyre Heros}}
\IfFontExistsTF{Consolas}{\setmonofont{Consolas}}{\setmonofont{Latin Modern Mono}}
\usepackage{amsmath,amssymb,mathtools}
\usepackage{geometry}
\usepackage{microtype}
\usepackage{enumitem}
\geometry{a4paper,margin=2.05cm}
\setlength{\parindent}{1.1em}
\setlength{\parskip}{0pt}
\newcommand{\foreign}[1]{#1}
\emergencystretch=220em
\sloppy
\hbadness=10000
\vbadness=10000
\hfuzz=5pt
\vfuzz=12pt
\raggedbottom
\begin{document}

\part*{Чест IV\hspace{1em}Библиографија}
\addcontentsline{toc}{part}{Чест IV: Библиографија}

\section*{Библиографија}
\addcontentsline{toc}{section}{Библиографија}

\begin{enumerate}[label={\small\arabic*.},leftmargin=2em,itemsep=0.3em]
\item О творјењу система форм тернарној биквадратичној формы.\\
\foreign{Sitz.\ Ber.\ d.\ Physikal.-mediz.\ Sozietät in Erlangen} \textbf{39} (1907), С.~176--179

\item О творјењу система форм тернарној биквадратичној формы.\\
\foreign{J.\ Reine Angew.\ Math.} \textbf{134} (1908), С.~23--90 и једна табела

\item К теорији инвариантов форм од $n$ вариаблов.\\
\foreign{J.\ Ber.\ d.\ DMV} \textbf{19} (1910), С.~101--104

\item К теорији инвариантов форм од $n$ вариаблов.\\
\foreign{J.\ Reine Angew.\ Math.} \textbf{139} (1911), С.~118--154

\item Рационалне функцијске поља.\\
\foreign{J.\ Ber.\ d.\ DMV} \textbf{22} (1913), С.~316--319

\item Поља и системы рационалных функциј.\\
\foreign{Math.\ Ann.} \textbf{76} (1915), С.~161--196

\item Теорем конечности инвариантов конечных груп.\\
\foreign{Math.\ Ann.} \textbf{77} (1916), С.~89--92

\item О цєлочислном рационалном прєдставјењу инвариантов система либовољно многых основных форм.\\
\foreign{Math.\ Ann.} \textbf{77} (1916), С.~93--102. (Глеј такоже Нр.~16)

\item Најобчејше области из целых трансцендентных чисел.\\
\foreign{Math.\ Ann.} \textbf{77} (1916), С.~103--128. (Глеј такоже Нр.~16)

\item Функционалне једначења изоморфного одображења.\\
\foreign{Math.\ Ann.} \textbf{77} (1916), С.~536--545

\item Једначења с предписаноју групоју.\\
\foreign{Math.\ Ann.} \textbf{78} (1918), С.~221--229. (Глеј такоже Нр.~16)

\item Инварианты либовољных диференциалных изразов.\\
\foreign{Nachr.\ v.\ d.\ Ges.\ d.\ Wiss.\ zu Göttingen} \textbf{1918}, С.~37--44

\item Инвариантне вариацијске проблемы.\\
\foreign{Nachr.\ v.\ d.\ Ges.\ d.\ Wiss.\ zu Göttingen} \textbf{1918}, С.~235--257

\item Аритметична теорија алгебраичных функциј једној вариаблы в јеј односу к другим теоријам и к теорији числовых пољ.\\
\foreign{J.\ Ber.\ d.\ DMV} \textbf{28} (1919), С.~182--203

\item Конечност система цєлочисловых инвариантов бинарных форм.\\
\foreign{Nachr.\ v.\ d.\ Ges.\ d.\ Wiss.\ zu Göttingen} \textbf{1919}, С.~138--156

\item К развоју в реды в теорији форм.\\
\foreign{Math.\ Ann.} \textbf{81} (1920), С.~25--30

\item Сполу с \foreign{W.\ Schmeidler}: модулы в некомутативных областјах, особливо из диференциалных и разликовых изразов.\\
\foreign{Math.\ Z.} \textbf{8} (1920), С.~1--35

\item О работє К.\ Хентзелт, падшего во војнє, к теорији елиминације.\\
\foreign{J.\ Ber.\ d.\ DMV} \textbf{30} (1921), С.~48

\item Теорија идеалов в колцовых областјах.\\
\foreign{Math.\ Ann.} \textbf{83} (1921), С.~24--66

\item Алгебраични критериј абсолутној иредуцибилности.\\
\foreign{Math.\ Ann.} \textbf{85} (1922), С.~26--33

\item Формално вариацијско рачуњање и диференциалне инварианты.\\
\foreign{Encyklopädie d.\ math.\ Wiss.} III,\ 3 (1922), С.~68--71 (в:
\foreign{R.\ Weitzenböck}, \textit{Диференциалне инварианты})

\item Обработање К.\ Хентзелт: к теорији полиномиалных идеалов и резултантов.\\
\foreign{Math.\ Ann.} \textbf{88} (1923), С.~53--79

\item Алгебраичне и диференциалне варианты.\\
\foreign{J.\ Ber.\ d.\ DMV} \textbf{32} (1923), С.~177--184

\item Теорија елиминације и обча теорија идеалов.\\
\foreign{Math.\ Ann.} \textbf{90} (1923), С.~229--261

\item Теорија елиминације и теорија идеалов.\\
\foreign{J.\ Ber.\ d.\ DMV} \textbf{33} (1924), С.~116--120

\item Абстрактно изграджење теорије идеалов в алгебраичном числовом пољу.\\
\foreign{J.\ Ber.\ d.\ DMV} \textbf{33} (1924), С.~102

\item Хилбертове функције в теорији идеалов.\\
\foreign{J.\ Ber.\ d.\ DMV} \textbf{34} (1925), С.~101

\item Карактери груп и теорија идеалов.\\
\foreign{J.\ Ber.\ d.\ DMV} \textbf{34} (1925), С.~144

\item Теорем конечности инвариантов конечных линеарных груп карактеристики $p$.\\
\foreign{Nachr.\ v.\ d.\ Ges.\ d.\ Wiss.\ zu Göttingen} \textbf{1926}, С.~28--35

\item Абстрактно изграджење теорије идеалов в алгебраичных числовых и функцијскых пољах.\\
\foreign{Math.\ Ann.} \textbf{96} (1927), С.~26--61

\item Теорем о дискриминанту за поредкы алгебраичного числового или функцијского поља.\\
\foreign{J.\ Reine Angew.\ Math.} \textbf{157} (1927), С.~82--104

\item Сполу с Р.\ Brauer: о минималных разпадных пољах иредуцибилных прєдставјениј.\\
\foreign{Sitz.\ Ber.\ d.\ Preuß.\ Akad.\ d.\ Wiss.} \textbf{1927}, С.~221--228

\item Хиперкомплексне величины и теорија прєдставјениј в аритметичном разумєњу.\\
\foreign{Atti Congresso Bologna} \textbf{2} (1928), С.~71--73

\item Хиперкомплексне величины и теорија прєдставјениј.\\
\foreign{Math.\ Z.} \textbf{30} (1929), С.~641--692

\item О максималных областјах из цєлочисловых функциј.\\
\foreign{Rec.\ Soc.\ Math.\ Moscou} \textbf{36} (1929), С.~65--72

\item Диференцирање идеалов и диферента.\\
\foreign{J.\ Ber.\ d.\ DMV} \textbf{39} (1930), С.~17

\item Нормална база в пољах без выше развєтвљења.\\
\foreign{J.\ Reine Angew.\ Math.} \textbf{167} (1932), С.~147--152

\item Сполу с Р.\ Brauer и Х.\ Hasse: доказ главного теорема в теорији алгебер.\\
\foreign{J.\ Reine Angew.\ Math.} \textbf{167} (1932), С.~399--404

\item Хиперкомплексне системы в својих односах к комутативној алгебри и к теорији чисел.\\
\foreign{Verhandl.\ Intern.\ Math.-Kongreß Zürich} \textbf{1} (1932), С.~189--194

\item Некомутативне алгебры.\\
\foreign{Math.\ Z.} \textbf{37} (1933), С.~514--541

\item Теорем о главном роду за релативно Галоисове числове поља.\\
\foreign{Math.\ Ann.} \textbf{108} (1933), С.~411--419

\item Разложене скрежене продукты и јих максималне поредкы.\\
\foreign{Actualités scientifiques et industrielles} \textbf{148} (1934) (15 С.)

\item Диференцирање идеалов и диферента.\\
\foreign{J.\ Reine Angew.\ Math.} \textbf{188} (1950), С.~1--21
\end{enumerate}

\vspace{1em}
\noindent\rule{\textwidth}{0.4pt}
\vspace{0.5em}

\noindent\textbf{Потрєбне и достаточне условја многократности к Нетеровому}\\
\textbf{фундаменталному теорему алгебраичных функциј: Х.\ Капферер}\\
\textbf{(с додатком, сполу с Е.\ Noether).}\\
\foreign{Math.\ Ann.} \textbf{97} (1927), С.~559--567

\newpage
\section*{Список краткых комуникациј и рецензиј Emmy Noether}
\addcontentsline{toc}{section}{Кратке комуникације и рецензије}

\subsection*{Кратке комуникације, објавјене в летописах Њемецкого математичного дружства}

\begin{itemize}[leftmargin=1.5em,itemsep=0.2em]
\item О инвариантах либовољных диференциалных изразов, 27, чест II, (1918), С.~28
\item Инвариантне вариацијске проблемы, 27, чест I, (1918), С.~47
\item Конечност цєлочисловых бинарных инвариантов, 26, чест II, (1919), С.~29
\item О цєлочисловых полиномах и степеновых редах, 28, чест II, (1919), С.~29--30
\item Теоремы разложења теорије модулов, 29, чест II, (1920), С.~54
\item Елементарне дєлитељи и обча теорија идеалов, 30, чест II, (1921), С.~32
\item Указање на формулу Липсцхитз, 30, чест II, (1921), С.~48
\item Аналог Хилбертового теорема в обчеј теорији идеалов, 32, чест II, (1923), С.~202
\item Групово-теоретичне замєткы к проблему аксиоматики, 33, чест I, (1924), С.~120
\item Галоисова теорија в либовољных пољах, 33, чест II, (1925), С.~120
\item Хилбертове функције в теорији идеалов, 34, чест II, (1926), С.~101
\item Диференциалне квоциенты идеалов и теорија развєтвљења, 38, чест II, (1929), С.~61
\end{itemize}

\subsection*{Рецензије књиг, објавјене в летописах Њемецкого математичного дружства}

\begin{itemize}[leftmargin=1.5em,itemsep=0.2em]
\item Е.\ Студы, \textit{Увод в теорију инвариантов линеарных трансформациј на основи векторного рачуњања.} Прва чест. (\foreign{Die Wissenschaft}, Bd.~71) \foreign{Braunschweig} 1923, \foreign{Vieweg}.\\
36, чест II, (1927), С.~71

\item Дєла Леополд Кронецкер, т.~4. Издал К.\ Хенсел. VI и 509 с. Leipzig и Берлин 1929, Б.\ Г.\ Teubner.\\
40, чест II, (1931), С.~112

\item Дєла Леополд Кронецкер, т.~5. Издал К.\ Хенсел. 527 с. Leipzig и Берлин 1930, Б.\ Г.\ Teubner.\\
41, чест I, (1932), С.~17

\item Дєла Леополд Кронецкер, т.~III.\,2. Издал К.\ Хенсел. 215 с. Leipzig и Берлин 1931, Б.\ Г.\ Teubner.\\
42, чест I, (1933), С.~38--39

\item Д.\ Е.\ Рутхерфорд, \textit{Модуларне инварианты.} \foreign{Cambridge Tracts No.~27.} Cambridge 1932, VII и 84 с.\\
43, чест II, (1934), С.~94--95
\end{itemize}

\newpage
\section*{Књигы, составјене с учестјем Emmy Noether}
\addcontentsline{toc}{section}{Књигы}

\begin{itemize}[leftmargin=1.5em,itemsep=0.5em]
\item Р.\ Дедекинд, \textit{Собране математичне дєла} I--III. Издали Р.\ Фрицке, Emmy Noether и \foreign{Øystein Ore}. \foreign{Vieweg}, \foreign{Braunschweig} 1930--1932\\
(т.~И 1930; т.~II 1931; т.~III 1932)

\item Г.\ Цантор, \textit{Листовање с Р.\ Дедекинд.} Издали Emmy Noether и \foreign{J.\ Cavaillès}, 60 страниц. \foreign{Actualités scient.\ et ind.} 581, Парис 1937
\end{itemize}

\vfill

\begin{center}
\rule{0.5\textwidth}{0.4pt}\\[1em]
\textit{Конец избраных математичных работ}\\[0.5em]
\textit{Emmy Noether}\\[0.5em]
\textit{(1882--1935)}
\end{center}

\end{document}
